\label{cap:fundamentacao}

Este capítulo apresenta os conceitos fundamentais que sustentam o desenvolvimento das \textit{fuzzing engines} para avaliação ética de sistemas de Inteligência Artificial (IA). A fundamentação teórica é organizada em quatro eixos principais: (i) princípios éticos de IA, (ii) métodos e frameworks para operacionalização ética, (iii) fundamentos de \textit{fuzzing} e testes automatizados aplicados a IA, e (iv) trabalhos correlatos mais diretamente relacionados ao escopo desta dissertação. A organização foi guiada pelo conjunto de 58 estudos primários selecionados no Mapeamento Sistemático da Literatura (MSL), apresentado no Capítulo~\ref{cap:rsl}.

%%%%%%%%%%%%%%%%%%%%%%%%%%%%%%%%%%%%%%%%%%%%%%%%%%%%%%%%%%%%%%%%%%%%%%%%%%%%%%%%%%%%%%%%%%%%%%%
\section{Princípios Éticos em Inteligência Artificial}
\label{sec:principios}

O ecossistema de diretrizes éticas em IA cresceu rapidamente desde 2016, consolidando conjuntos de princípios que aparecem de forma recorrente em documentos internacionais \cite{Jobin2019, UNESCO2021, EUAIAct2024} e \textit{Australia’s AI Ethics} Principles \footnote{\url{https://www.industry.gov.au/publications/australias-artificial-intelligence-ethics-principles/australias-ai-ethics-principles}}. No entanto, os estudos analisados no MSL demonstram que, embora muitos princípios sejam discutidos globalmente, apenas alguns são consistentemente operacionalizados em contextos de Engenharia de Software. Entre os princípios identificados nos 58 estudos, três se destacam por sua prevalência e profundidade de aplicação técnica: \textit{Fairness}, \textit{Accountability} e \textit{Transparency}. Estes são analisados com maior detalhe nas subseções seguintes por serem diretamente implementados nesta pesquisa.

\subsection{Fairness}
\label{sec:fairness}

O princípio de \textit{Fairness} emergiu como o mais recorrente no corpus analisado, presente em 48 dos 58 estudos. A literatura enfatiza o risco de discriminação algorítmica decorrente de dados enviesados \cite{S003, S007, S034}, modelos historicamente treinados com atributos sensíveis \cite{S015, S022} e comportamentos diferenciados entre grupos protegidos \cite{S030, S034}. Entre as abordagens identificadas, destacam-se:

\begin{itemize}
    \item métricas estatísticas de equidade, como \textit{demographic parity}, \textit{equalized odds} e \textit{predictive parity} \cite{S003, S004, S034};
    \item métodos de mitigação de viés em diferentes estágios do pipeline, como pré-processamento \cite{S022, S026}, \textit{in-processing} \cite{S009, S015} e pós-processamento \cite{S004, S034};
    \item auditorias automatizadas aplicadas a modelos de classificação e regressão, principalmente nos domínios de saúde, seleção de pessoal e crédito \cite{S003, S012, S038}.
\end{itemize}

Apesar dos avanços, os estudos destacam limitações importantes, como a falta de consenso sobre qual métrica usar em cada cenário \cite{S003, S004} e a dificuldade em analisar \textit{trade-offs} entre equidade e acurácia \cite{S015, S034}.

\subsection{Accountability}
\label{sec:accountability}

O princípio de \textit{Accountability} foi identificado em 32 dos 58 estudos. Esse princípio refere-se à ecessidade de atribuir responsabilidade por decisões algorítmicas, garantindo rastreabilidade, auditoria e governança ao longo do ciclo de vida dos sistemas de IA \cite{S001, S014, S046}. Enquanto parte da literatura enfatiza mecanismos organizacionais de supervisão \cite{S014, S046}, outras contribuições propõem:

\begin{itemize}
    \item mecanismos de \textit{logging} ético e auditoria de decisões \cite{S005, S046};
    \item padrões arquiteturais que garantem verificabilidade \cite{S014, S043};
    \item frameworks de governança baseados em responsabilidades distribuídas
    \cite{S025, S038}.
\end{itemize}

O MSL mostra que, embora existam processos e frameworks conceituais, há escassez de ferramentas automatizadas que verificam responsabilidade técnica durante testes ou execução, lacuna endereçada diretamente por esta pesquisa.

\subsection{Transparency}
\label{sec:transparency}

O princípio de \textit{Transparency} apareceu em 29 estudos e frequentemente se sobrepõe a \textit{Explainability} \cite{S006, S012}. No corpus analisado, este princípio engloba:

\begin{itemize}
    \item explicações locais e globais (LIME, SHAP, Grad-CAM) \cite{S008, S027};
    \item modelos intrinsecamente interpretáveis \cite{S009};
    \item documentação estruturada de modelos, como \textit{model cards} \cite{S046};
    \item requisitos de transparência previstos em diretrizes internacionais
    \cite{UNESCO2021, EUAIAct2024}.
\end{itemize}

Os estudos mais recentes destacam a necessidade de medir transparência de maneira sistemática, incluindo avaliação da clareza das explicações e consistência entre execuções \cite{S008, S012}, além de incorporar dimensões além da interpretabilidade técnica, como comunicação com usuários leigos.

%%%%%%%%%%%%%%%%%%%%%%%%%%%%%%%%%%%%%%%%%%%%%%%%%%%%%%%%%%%%%%%%%%%%%%%%%%%%%%%%%%%%%%%%%%%%%%%
\section{Abordagens e Métodos para Operacionalização Ética}
\label{sec:operacionalizacao}

Embora princípios éticos sejam amplamente discutidos, sua operacionalização ainda é incipiente, um ponto enfatizado repetidamente pelos estudos analisados \cite{S012, S014, S030}. Os 58 estudos foram classificados em três categorias principais.

\subsection{Métodos e Frameworks Estruturados}

Diversos estudos propõem métodos formais para incorporar considerações éticas ao ciclo de vida de sistemas de IA. Destaques incluem:

\begin{itemize}
    \item o método ECCOLA \cite{S012, S013}, que fornece checklists e diretrizes práticas para suportar     equipes de desenvolvimento na implementação de IA ética;
    \item frameworks de \textit{ethics by design} \cite{S014, S025}, que integram princípios éticos desde a fase de requisitos;
    \item processos completos alinhados ao SDLC tradicional \cite{S006, S046}.
\end{itemize}

Contudo, a maioria dos métodos estruturados carece de ferramentas automatizadas e depende de análise manual, o que reduz aplicabilidade industrial.

\subsection{Ferramentas e Bibliotecas Técnicas}

A análise também identificou ferramentas específicas voltadas para operacionalização prática de princípios éticos, especialmente:

\begin{itemize}
    \item bibliotecas de mitigação de viés e avaliação de equidade (\textit{AIF360}, \textit{Fairlearn}) frequentemente citadas pelos estudos técnicos \cite{S003, S009, S015};
    \item ferramentas de explicabilidade (\textit{LIME}, \textit{SHAP}, \textit{ELI5}) aplicadas em diversos casos práticos \cite{S008, S027};
    \item protótipos de auditoria automatizada, ainda majoritariamente experimentais \cite{S005, S046, S038}.
\end{itemize}

No entanto, nenhum dos estudos apresenta ferramentas integradas a pipelines de CI/CD ou IDEs, explicitando a lacuna mencionada em \cite{S012, S021, S030}.

\subsection{Integração ao Ciclo de Vida do Desenvolvimento}

A análise de fases do SDLC (RQ2 do MSL) mostrou que a maioria das abordagens se concentra em:

\begin{itemize}
    \item \textbf{Requisitos (55,2\% dos estudos)}: identificação e elicitação de necessidades éticas     \cite{S001, S006, S014};
    \item \textbf{Design (50,0\%)}: decisões arquiteturais e padrões para suportar responsabilidade e     transparência \cite{S006, S014, S043};
    \item \textbf{Implementação (56,9\%)}: aplicações técnicas diretas, como mitigação de viés,     explicabilidade e logging \cite{S003, S009, S027};
    \item \textbf{Testes (37,9\%)}: foco restrito a avaliação de \textit{fairness} e explicabilidade, quase sempre em cenários experimentais \cite{S003, S004, S022}.
\end{itemize}

As fases posteriores  \textit{Deployment}, \textit{Operation}, \textit{Monitoring}, \textit{Governance} são muito menos representadas, apesar de críticas para manutenção ética contínua.

%%%%%%%%%%%%%%%%%%%%%%%%%%%%%%%%%%%%%%%%%%%%%%%%%%%%%%%%%%%%%%%%%%%%%%%%%%%%%%%%%%%%%%%%%%%%%%%
\section{Fundamentos de \textit{Fuzzing} e Testes Automatizados}
\label{sec:fuzzing}

\textit{Fuzzing} é uma técnica consolidada em testes de software para detecção automática de falhas por meio da geração massiva de entradas \cite{Zhang2019}. Embora amplamente usada em segurança, seu uso para avaliação ética ainda é emergente.

\subsection{Fuzzing Tradicional}

Os elementos fundamentais do \textit{fuzzing} incluem:

\begin{itemize}
    \item \textbf{geradores de entradas} (mutação ou geração sintética);
    \item \textbf{oráculos} capazes de determinar se a saída representa uma falha;
    \item \textbf{instrumentação} para coletar métricas e caminhos;
    \item \textbf{feedback loops} para priorização de casos.
\end{itemize}

Técnicas como \textit{coverage-guided fuzzing} e \textit{model-based fuzzing} têm sido aplicadas com sucesso em sistemas críticos, mas apenas recentemente migraram para IA.

\subsection{Fuzzing para IA}

Zhang et al. \cite{Zhang2019} demonstram que \textit{fuzzing} pode ser adaptado para detectar:

\begin{itemize}
    \item vieses estatísticos;
    \item inconsistências explicativas;
    \item comportamentos não transparentes;
    \item efeitos adversos de decisões algorítmicas.
\end{itemize}

Entretanto, a literatura identificada pelo MSL revela que ainda não existem frameworks completos ou engines modulares para avaliação ética, reforçando a contribuição desta dissertação.

%%%%%%%%%%%%%%%%%%%%%%%%%%%%%%%%%%%%%%%%%%%%%%%%%%%%%%%%%%%%%%%%%%%%%%%%%%%%%%%%%%%%%%%%%%%%%%%
\section{Trabalhos Correlatos}
\label{sec:correlatos}

Esta seção apresenta uma análise aprofundada dos estudos mais relacionados ao escopo desta pesquisa, considerando (i) métodos estruturados para operacionalização ética, (ii) ferramentas práticas para avaliação ou mitigação de riscos éticos, e (iii) técnicas automatizadas próximas à lógica de \textit{fuzzing}. Os trabalhos discutidos aqui representam o estado da arte mais diretamente conectado ao desenvolvimento das \textit{fuzzing engines} propostas nesta dissertação, cobrindo o contexto, método, resultados, limitações e relação com esta pesquisa.

%%%%%%%%%%%%%%%%%%%%%%%%%%%%%%%%%%%%%%%%%%%%%%%%%%%%%%%%%%%%%%%%%%%%%%%%%%%%%%%%%%%%%%%%%%%%%%%
\subsection{Métodos Estruturados Relacionados}

\paragraph{\textbf{Método ECCOLA (S012, S013).}}
O método ECCOLA, introduzido em \cite{S012} e posteriormente expandido em \cite{S013}, propõe um conjunto estruturado de cartões (\textit{cards}) contendo diretrizes éticas aplicáveis ao desenvolvimento de sistemas de IA. Os estudos utilizam sessões participativas com profissionais para validar a utilidade dos cartões na elicitação de requisitos éticos. Os resultados mostram que os cartões ajudam equipes a identificar riscos éticos durante requisitos e design, mas o método depende de processos manuais, sem mecanismos de verificação automática ou integração ao SDLC. Embora seja um dos métodos mais consolidados para operacionalização conceitual, carece de suporte ferramental automatizado -- lacuna diretamente abordada pelas \textit{fuzzing engines} deste trabalho.

\paragraph{\textbf{Ethics-Based Auditing (S014).}}
O estudo \cite{S014} propõe uma estrutura analítica para auditorias éticas de sistemas automatizados de decisão, aplicando entrevistas, análise documental e mapeamento de \textit{intervention points}. O método permite identificar riscos éticos ao longo do ciclo de vida da solução e orientar intervenções organizacionais. Os resultados destacam desafios como falta de rastreabilidade e limitações técnicas na coleta de evidências. Embora forneça diretrizes organizacionais robustas, o estudo não fornece mecanismos automatizados de auditoria técnica, reforçando a necessidade de técnicas como \textit{fuzzing} ético para avaliação sistemática.

\paragraph{\textbf{Value-Sensitive and User-Centred Design (S025, S041).}}
Os estudos \cite{S025, S041} aplicam princípios de \textit{Value-Sensitive Design} (VSD) e design centrado no usuário à incorporação de valores éticos em software. Empregam métodos qualitativos, como entrevistas e oficinas, para identificar valores de usuários marginalizados e traduzi-los em requisitos. Embora contribuam com metodologias interpretativas essenciais, as abordagens permanecem inteiramente qualitativas e desconectadas de verificações técnicas automatizadas, reforçando o papel complementar das \textit{fuzzing engines} em garantir conformidade durante testes.

\paragraph{\textbf{Ethically Aligned Design e FATES (S046, S044).}}
Os estudos \cite{S046, S044} propõem integrar princípios éticos (como os da IEEE EAD) em processos ágil e CRISP-DM. Eles revisam etapas de desenvolvimento, propondo adaptações metodológicas (artefatos adicionais, práticas de documentação, critérios de aceitação éticos). Contudo, os autores relatam baixa padronização e dependência de especialistas, o que limita escalabilidade. As abordagens não incluem mecanismos de verificação automatizada — lacuna crítica abordada pelo \textit{fuzzing} proposto nesta dissertação.

%%%%%%%%%%%%%%%%%%%%%%%%%%%%%%%%%%%%%%%%%%%%%%%%%%%%%%%%%%%%%%%%%%%%%%%%%%%%%%%%%%%%%%%%%%%%%%%
\subsection{Ferramentas Práticas Relacionadas}

\paragraph{\textbf{AIF360 (S003) e Fairlearn (S009, S015).}}
Os estudos \cite{S003, S009, S015} avaliam bibliotecas amplamente utilizadas para mitigação de viés. AIF360 fornece métricas de \textit{fairness}, algoritmos de mitigação e exemplos aplicados, enquanto Fairlearn enfoca mitigação por restrições e análise de desempenho por grupo. Em todos os estudos, as ferramentas são testadas em \textit{datasets} de benchmark, e os autores destacam que mitigação tende a reduzir acurácia. Apesar de úteis, tais bibliotecas não oferecem geração de casos de teste automatizados nem avaliação exploratória -- requisitos centrais para \textit{fuzzing} ético.

\paragraph{\textbf{Frameworks de Explicabilidade (S006, S008, S027).}}
Os estudos \cite{S006, S008, S027} exploram técnicas de explicabilidade, como LIME, SHAP e Grad-CAM. Os métodos são avaliados quanto à clareza de explicações, estabilidade entre execuções e fidelidade. Embora esses estudos aprofundem interpretações locais e globais, os autores relatam limitações importantes: instabilidade de explicações \cite{S008}, falta de padronização e dificuldade de uso por não especialistas. Essas limitações reforçam a necessidade de testes automatizados capazes de avaliar transparência de forma sistemática, como proposto pelas \textit{fuzzing engines} desta pesquisa.

\paragraph{\textbf{Model Cards e Transparência Documentada (S046).}}
O estudo \cite{S046} revisita e expande o uso de \textit{Model Cards} para documentação ética de sistemas de IA. Os autores aplicam estudos de caso e identificam que documentação estruturada melhora comunicação e governança, mas depende de autores experientes e é vulnerável a omissões. Como o trabalho não inclui mecanismos de verificação, ele demonstra a importância de combinar documentação com testes automatizados, abordagem adotada neste trabalho.

%%%%%%%%%%%%%%%%%%%%%%%%%%%%%%%%%%%%%%%%%%%%%%%%%%%%%%%%%%%%%%%%%%%%%%%%%%%%%%%%%%%%%%%%%%%%%%%
\subsection{Técnicas Automatizadas Próximas ao Fuzzing}

\paragraph{\textbf{Fuzzing Ético e Testes Baseados em Perturbações (Zhang et al., 2019).}}
O trabalho de Zhang et al. \cite{Zhang2019} é o mais próximo ao escopo desta dissertação. Os autores adaptam \textit{coverage-guided fuzzing} para testar modelos de IA, propondo geração sistemática de entradas que expõem vieses e inconsistências. O estudo demonstra viabilidade da técnica, mas aplica apenas um protótipo experimental, sem modularização, integração ao SDLC ou aplicação a múltiplos princípios éticos. O presente trabalho avança significativamente ao desenvolver três engines completas, modulares e integráveis, com foco explícito em \textit{Fairness}, \textit{Transparency} e \textit{Accountability}.

\paragraph{\textbf{Auditorias Automatizadas Experimentais (S005, S046).}}
Os estudos \cite{S005, S046} apresentam iniciativas iniciais de auditorias parciais de sistemas de IA. Em \cite{S005}, são propostos métodos para avaliar sustentabilidade e impacto sistêmico em modelos, utilizando métricas computacionais e indicadores de impacto. Em \cite{S046}, auditorias são acopladas a pipelines experimentais. Em ambos os casos, os autores apontam limitações como baixa cobertura, ausência de testes exploratórios e falta de escalabilidade, lacunas justamente endereçadas pela aplicação de \textit{fuzzing} ético nesta dissertação.

%%%%%%%%%%%%%%%%%%%%%%%%%%%%%%%%%%%%%%%%%%%%%%%%%%%%%%%%%%%%%%%%%%%%%%%%%%%%%%%%%%%%%%%%%%%%%%%
A análise dos estudos correlatos revela que, embora existam métodos, frameworks e ferramentas úteis para operacionalização ética, nenhuma abordagem atual oferece \textbf{testes automatizados, exploratórios e sistemáticos} para avaliação ética de modelos de IA já implementados. Métodos como ECCOLA estruturam processos conceituais, ferramentas como AIF360 executam mitigação pontual, e frameworks de explicabilidade geram interpretações isoladas, porém todos carecem de automatização contínua, integração ao SDLC e mecanismos de verificação massiva. Assim, a literatura reforça a originalidade e necessidade da proposta desta dissertação, que avança o estado da arte ao desenvolver \textit{fuzzing engines} dedicadas para \textit{Fairness}, \textit{Transparency} e \textit{Accountability}.

%%%%%%%%%%%%%%%%%%%%%%%%%%%%%%%%%%%%%%%%%%%%%%%%%%%%%%%%%%%%%%%%%%%%%%%%%%%%%%%%%%%%%%%%%%%%%%%
\section{Síntese do Capítulo}

O MSL revelou que princípios éticos são amplamente discutidos, mas raramente operacionalizados de forma prática e automatizada. Entre todos os princípios identificados, \textit{Fairness}, \textit{Accountability} e \textit{Transparency} apresentam maior maturidade técnica, sendo suportados por métricas, padrões arquiteturais e ferramentas emergentes. No entanto, as abordagens atuais permanecem limitadas por: i) dependência de processos manuais; ii) foco restrito às fases iniciais do SDLC; iii) baixa integração a pipelines de desenvolvimento; e iv) ausência de verificação automatizada, especialmente na fase de testes. A utilização de técnicas de \textit{fuzzing} surge como oportunidade para preencher lacunas identificadas no MSL, permitindo avaliação sistemática, escalável e repetível de riscos éticos associados a IA. Essa base conceitual fundamenta diretamente a proposta apresentada no Capítulo~\ref{cap:proposta}.